% Part: lambda-calculus
% Chapter: introduction
% Section: church-rosser

\documentclass[../../../include/open-logic-section]{subfiles}

\begin{document}

\olfileid[id]{lam}{int}{cr}
\olsection{Sifat Church--Rosser}

\begin{thm}
\ollabel{thm:church-rosser}
Misalkan $M$, $N_1$, dan $N_2$ adalah suku sedemikian sehingga $M \red N_1$
dan $M \red N_2$. Maka ada suatu suku $P$ sedemikian sehingga $N_1 \red P$
dan $N_2 \red P$.
\end{thm}

\begin{cor}
Andaikan $M$ dapat direduksi ke bentuk normal. Maka bentuk normal ini tunggal.
\end{cor}

\begin{proof}
Jika $M \red N_1$ dan $M \red N_2$, berdasarkan teorema sebelumnya terdapat
suatu suku~$P$ sedemikian sehingga $N_1$ dan $N_2$ keduanya tereduksi
menjadi~$P$. Jika $N_1$ dan~$N_2$ keduanya berbentuk normal, hal ini hanya
mungkin jika $N_1 \ident P \ident N_2$.
\end{proof}

Terakhir, kita mengatakan bahwa dua suku $M$ dan~$N$ \emph{$\beta$-ekuivalen},
atau cukup \emph{ekuivalen}, jika keduanya tereduksi menjadi suatu suku yang
sama; dengan kata lain, jika terdapat $P$ sedemikian sehingga $M \red P$ dan
$N \red P$. Hal ini ditulis $M \equal[\beta] N$. Dengan menggunakan
\olref{thm:church-rosser}, Anda dapat memeriksa bahwa $\equal[\beta]$ adalah
relasi ekuivalensi dengan sifat tambahan bahwa, untuk setiap $M$ dan~$N$, jika
$M \red N$ atau $N \red M$, maka $M \equal[\beta] N$. (Bahkan, dapat
ditunjukkan bahwa $\equal[\beta]$ merupakan relasi ekuivalensi \emph{terkecil}
yang mempunyai sifat tersebut.)

\end{document}

